\thispagestyle{plain}

\ifthenelse{\equal{\getLanguage}{portuguese}}{%
    \pdfbookmark[1]{Resumo}{resumo}
    \chapter*{Resumo}
}{%
    \pdfbookmark[1]{Abstract}{abstract}
    \chapter*{Tóm tắt báo cáo}
}

Báo cáo Bài tập lớn môn \textbf{Các hệ thống Cơ sở dữ liệu} thực hiện nghiên cứu, phân tích và triển khai toàn diện \textbf{Hệ cơ sở dữ liệu quản lý một giải bóng đá chuyên nghiệp} (Đề tài 13). Hệ thống được xây dựng nhằm đáp ứng các yêu cầu nghiệp vụ thực tế của một mùa giải bóng đá: quản lý câu lạc bộ, cầu thủ, sân vận động, đội ngũ trọng tài, hợp đồng thi đấu và cho mượn cầu thủ; thiết lập lịch thi đấu vòng tròn hai lượt ($10$ vòng, $30$ trận); đăng ký đội hình tối đa 18 cầu thủ mỗi trận; theo dõi chi tiết các sự kiện diễn biến trên sân (bàn thắng, thẻ phạt) và tính toán bảng xếp hạng tự động. 

Nội dung báo cáo bao gồm: (1) Phân tích yêu cầu và thiết kế Lược đồ thực thể liên kết chi tiết (E-R); (2) Ánh xạ sang Lược đồ quan hệ và chứng minh giải tích toàn bộ $11$ quan hệ đạt \textbf{Dạng chuẩn BCNF (Boyce-Codd Normal Form)}; (3) Cài đặt CSDL trên hệ quản trị với đầy đủ ràng buộc toàn vẹn, ít nhất $5$ Transaction và $5$ Trigger nghiệp vụ phức tạp; (4) Phân tích và thiết kế \textbf{Cơ sở dữ liệu phân tán (CSDLPT)} theo mô hình đa chi nhánh gồm phân mảnh ngang, phân mảnh dọc, phân mảnh dẫn xuất và nhân bản dữ liệu; (5) Xây dựng ứng dụng giao diện trực quan kết nối CSDL phục vụ tra cứu, cập nhật dữ liệu, tính lương trọng tài và thống kê doanh thu bán vé.

\vspace{1cm}

\noindent\textbf{Từ khóa:} Cơ sở dữ liệu, Lược đồ E-R, Chuẩn hóa BCNF, CSDL Phân tán, Phân mảnh dữ liệu, Quản lý giải bóng đá, Trigger, Transaction.
